% TOP_PROBLEM: {"id": 5500017, "title": "Visibility Graph Recognition", "category_id": 3, "category": {"id": 3, "name": "graph_theory", "display_name": "Graph Theory", "description": "Problems involving graphs, networks, and their properties.", "slug": "graph-theory", "order_index": 3, "created_at": "2026-07-31T15:26:25.671Z"}, "status": "open", "problem_number": "AMR-054-0017", "set_id": 13}
% FIRST_POSTED: 2026-10-01
% ENTRY_KIND: statement_only
% UnsolvedMath ID 5500017; source catalogue code AMR-054-0017
% !TeX program = lualatex
% Definition and sources only; this file is not an LLM solution attempt.
\documentclass[11pt,a4paper]{article}
\usepackage[margin=25mm]{geometry}
\usepackage{fontspec}
\setmainfont{lmroman10-regular.otf}[BoldFont=lmroman10-bold.otf,
 ItalicFont=lmroman10-italic.otf,BoldItalicFont=lmroman10-bolditalic.otf]
\usepackage{amsmath,amssymb,mathtools}
\usepackage{unicode-math}
\setmathfont{latinmodern-math.otf}
\AtBeginDocument{\let\setminus\smallsetminus}
\usepackage{xurl,hyperref}
\hypersetup{colorlinks=true,urlcolor=blue}
\setcounter{secnumdepth}{0}
\setlength{\parindent}{0pt}
\setlength{\parskip}{5pt}
\setlength{\emergencystretch}{3em}
\begin{document}
\section*{Visibility Graph Recognition}\label{5500017}
{\small UnsolvedMath ID 5500017 · AMR-054-0017 · Graph Theory · AMR Open Problem Lists}
\subsection{Definitions and mathematical statement}
The input is a finite simple graph $G$ on labelled vertices
$1,\ldots,n$, together with a specified Hamiltonian cycle
$C=(v_1,\ldots,v_n,v_1)$ in $G$. A realization consists of distinct
points $p_1,\ldots,p_n\in\mathbb R^2$ such that the segments following
$C$ form the boundary of a simple polygon $P$.

Two polygon vertices are visible when the entire closed segment between
them lies in $P$, including its boundary. Require the exact equivalence
\[
            ij\in E(G)\quad\Longleftrightarrow\quad[p_i,p_j]\subset P
            \qquad(i\ne j).
\]
Thus every edge and every nonedge constrains the realization; merely
drawing all graph edges inside some polygon is insufficient.

Can existence of such a realization be decided in time polynomial in
the size of the graph input? More generally, determine the complexity
of this recognition problem. The cycle supplies the intended boundary
order in advance, but no coordinates. The polygon has no holes, and
no extra vertices may be inserted. The decision problem does not
automatically require a rational-coordinate witness of polynomial
bit length; that would be a separate and stronger certificate claim.
\subsection{Short English statement}
Given a graph and its proposed polygon boundary cycle, decide whether
the graph records exactly which pairs of vertices can see one another
inside a simple polygon.
\subsection{Sources}
\begin{itemize}
\item[S1] ULAM.AI and the UnsolvedMath Contributors, supplied problems.json, record 5500017 (AMR-054-0017), AMR Open Problem Lists. Catalogue identity and source transcription; CC BY 4.0.
\url{https://huggingface.co/datasets/ulamai/UnsolvedMath}.
\item[S2] Open Problems Project - Computational Geometry, Problem 17. Original source indexed by AMR-054-0017.
\url{https://topp.openproblem.net/p17}.
\item[S3] Schaefer, Cardinal and Miltzow, The Existential Theory of the Reals as a Complexity Class: A Compendium, polygon visibility recognition.
\url{https://arxiv.org/abs/2407.18006}.
\end{itemize}
\end{document}
